\section{Method and benchmark}
\label{sec:method}

\subsection{Geometric oracle and identifiability}
\label{sec:oracle}

\textbf{Structures and labels.} A crystal structure is a triple $s = (L, X, Z)$: a lattice
$L \in \mathbb{R}^{3\times3}$, fractional coordinates $X = \{x_i\}_{i=1}^{N} \subset [0,1)^3$, and species
$Z = \{z_i\}$. The label map is $y(s) = \spg(s)$, the crystal system returned by a deterministic symmetry
algorithm (spglib) at a fixed tolerance $\tau$, taking values in the seven crystal systems. Because the
label is computed rather than annotated, label correctness is not a competing explanation for model error;
Section~\ref{sec:benchmark} describes the tolerance sweep that certifies it.

\textbf{Renders.} The frozen protocol renders a $2\times2\times2$ tiling of the conventional cell through
five orthographic cameras $\mathcal{C} = \{c_v\}_{v=1}^{5}$ (three principal-axis and two oblique views) at
$768$ px, giving images $I_v = \mathcal{R}(s, c_v)$. The tiling is a legibility choice inherited from the
render convention and is not information-free; Section~\ref{sec:conventions} measures what it costs. Each camera is a known unit direction $d_v$ with orthographic projection
$\Pi_v : \mathbb{R}^3 \to \mathbb{R}^2$, $\Pi_v(p) = P_v\, p$ for a fixed $2\times3$ matrix $P_v$ with
orthonormal rows perpendicular to $d_v$. The task given to a model is: from $(I_1,\dots,I_5)$, name the
crystal system.

\textbf{The oracle.} The oracle answers the same question from the renders' own geometry instead of from
pixels. It forward-projects the ground-truth atom positions of the conventional cell through each frozen
camera, \emph{discards} the correspondence between views, re-solves it by ray intersection with cross-view
verification, and applies $\spg$ to the reconstruction (Algorithm~\ref{alg:oracle}). It is deterministic
and CPU-only. Two properties of that definition are load-bearing and we state them here rather than in the
limitations. First, the oracle assumes \emph{ideal extraction}: every atom's centroid is available in every
view, including atoms an image hides behind another, so $R_1$ bounds identifiability from the projection
geometry and not from the pixels. A visibility-conditioned variant that removes occluded detections is a
separate measurement and is reported in Section~\ref{sec:conventions}. Second, the oracle inverts the
conventional-cell projections rather than the tiled ones the renderer draws; supplying it the tiled geometry
instead is one of the interventions of Section~\ref{sec:conventions} and costs it twenty-two structures.

\begin{algorithm}[t]
\caption{Geometric oracle $O(s, \mathcal{C})$}
\label{alg:oracle}
\begin{algorithmic}[1]
\For{each view $v$}
  \State $Q_v \gets \{(\Pi_v(L x_i),\, z_i) : i = 1,\dots,N\}$
  \Comment{projected, species-labelled 2D points; cross-view identity discarded}
\EndFor
\For{each pair of views $(v, w)$ and each same-species pair $(q, q') \in Q_v \times Q_w$}
  \State $\hat{p} \gets$ intersection of the back-projected rays $\Pi_v^{-1}(q)$ and $\Pi_w^{-1}(q')$
  \State accept $\hat{p}$ iff its projection lies within tolerance of a same-species point in every
         remaining view
\EndFor
\State \Return $\spg(\text{verified reconstruction})$
\end{algorithmic}
\end{algorithm}

\begin{definition}[Identifiability ceiling]
\label{def:ceiling}
For an evaluation sample $\mathcal{S}$, the ceiling is
\[
R_1 = \frac{1}{|\mathcal{S}|}\sum_{s \in \mathcal{S}} \mathbf{1}\!\left[\,O(s,\mathcal{C}) = y(s)\,\right],
\]
the fraction of structures whose label the frozen renders determine for an ideal geometric reader.
\end{definition}

\begin{proposition}[Exact recovery]
\label{prop:ident}
Let $s$ be a structure rendered under known orthographic cameras. Suppose (i) at least two of the camera
directions are linearly independent, and (ii) the cross-view verification of
Algorithm~\ref{alg:oracle} admits a unique consistent correspondence, i.e.\ no two distinct atoms produce
verified ray intersections within the matching tolerance of one another. Then the oracle recovers every
atom position exactly, and $O(s, \mathcal{C}) = \spg(s) = y(s)$.
\end{proposition}

A proof sketch is in Appendix~\ref{app:proof}; the argument is elementary (two orthographic images along
independent directions determine a point as the intersection of two affine lines, and verification against
the remaining views eliminates spurious matches outside the coincidence set). The proposition is useful for
what its failure set says. Condition (i) holds by construction for the frozen camera set. Condition
(ii) fails under projective coincidence, when distinct atoms project onto near-identical pixels, and this
is exactly where the measured ceiling falls short of $1$: coincidence grows with atoms per cell, which
predicts the cell-size dependence observed at scale-up (Section~\ref{sec:controls}). A separate quantity
governs what an extractor reading pixels could recover: across the frozen five-view set, the fraction of
atoms with fewer than two clear views is 0.26\% (original sample) and 0.87\% (replication sample), so the
protocol's occlusion leaves almost every atom triangulable in principle (Section~\ref{sec:conventions}).

\textbf{Why a ceiling below $1.0$ is the point.} A design that supplied the source coordinates as text
would have a ceiling of exactly $1.0$ by construction and would measure nothing about the images. A ceiling
below $1.0$ says the protocol withholds information, and Proposition~\ref{prop:ident} locates where.

\subsection{Attribution ladder}
\label{sec:ladder}

\begin{definition}[Rungs and shares]
\label{def:ladder}
For a model $m$ evaluated on the same structures, let
$R_0 = \tfrac{1}{n}\sum \mathbf{1}[\spg(s) = y(s)]$ (the symmetry algorithm on ground-truth positions),
$R_1$ the oracle ceiling of Definition~\ref{def:ceiling},
$R_3(m)$ the accuracy of $m$ given ground-truth fractional coordinates, species and cell parameters as
text, and
$R_4(m)$ the accuracy of $m$ on the pixel renders. Define the \emph{perception component}
$P(m) = R_3(m) - R_4(m)$, the \emph{symbolic residual} $S(m) = R_1 - R_3(m)$, and the \emph{perception
share} $P(m) / (R_1 - R_4(m))$. $R_2$ denotes the oracle run on a learned detector's output.
\end{definition}

Three conditions make the shares interpretable, and we state them because each is checkable: the rungs are
scored on the same named structures with all comparisons paired; $R_3$ and $R_4$ use the same decode budget
($K{=}3$ everywhere the shares are computed); and the symbolic residual \emph{bounds} rather than isolates
symmetry reasoning, since a residual against a deterministic solver can include failure modes other than
symmetry (we test this directly below). Table~\ref{tab:rungs} gives the ladder.

\begin{table}[t]
\centering
\caption{Attribution ladder, original evaluation sample, $n=210$, $K{=}3$. The two model columns are the
strongest scored model and one mid-roster model, named rather than aggregated; $R_3$ and $R_4$ within
a column are the same model under two conditions, so their interval is that model's perception component
(Definition~\ref{def:ladder}). $R_2$ is unscored: its extractor fails to triangulate on 19.0\% and 22.4\%
of structures, over the pre-registered 5\% gate. The fine-tuned arm was never run under the coordinate
condition, so it carries no within-model interval and is not a column.}
\label{tab:rungs}
\small
\setlength{\tabcolsep}{4.5pt}
\begin{tabular}{p{0.40\linewidth}p{0.22\linewidth}rr}
\toprule
rung & what it removes & \small gemini-3.6-flash & \small gpt-4.1-mini \\
\midrule
$R_0$ algorithm on ground-truth positions & nothing & \multicolumn{2}{c}{1.0000 (definitional)} \\
$R_1$ oracle: camera inversion, then algorithm & the protocol's losses & \multicolumn{2}{c}{0.9524} \\
$R_3$ exact geometry as text ($K{=}3$) & the model's perception & 0.8524 & 0.4143 \\
$R_4$ pixels ($K{=}3$) & nothing further & 0.7333 & 0.3667 \\
\midrule
perception share $(R_3-R_4)/(R_1-R_4)$ & --- & 0.5436 & 0.0813 \\
$R_2$ oracle on a learned detector's output & ideal extraction & \multicolumn{2}{c}{unscored} \\
\bottomrule
\end{tabular}
\end{table}

\textbf{The top rung is definitional, and that is load-bearing.} The label \emph{is} $\spg$ applied to the
structure, so $R_0 = 1.0000$ exactly, and the $R_0 \to R_1$ interval is what camera inversion and
correspondence re-solution cost. Treating $R_0$ as approximately equal to $R_1$ would discard that
interval.

\subsection{Benchmark construction}
\label{sec:benchmark}

\textbf{Certified labels.} Each structure's crystal system, Bravais lattice, point group and space group
are derived from its coordinates by $\spg$ at the canonical tolerance, with a tolerance sweep that
quarantines any structure whose space group is not stable across it. On the kept set, label agreement with
the source databases is $220/220$ (one-sided exact 95\% lower bound $0.9865$).

\textbf{Composition exclusion.} No chemical composition in any evaluation set appears in training. Without
this, symmetry is partly predictable from composition alone, a published task \citep{cryspnet2003}, and a
model could score well without reading the image.

\textbf{A label ladder gives a difficulty axis.} We score four labels of increasing granularity for the
model-free arms. The oracle is flat across granularity because a
reconstruction is either right or wrong and all four labels follow from it; the cell-metric baseline falls
from 0.8952 at crystal system to 0.6810 at space group, because point group and space group depend on atom
positions that lattice parameters do not encode. The identifiability claim is therefore not specific to the
seven-way label.
